\scanpage{176}
\mspage{90}
% L01–02
\noindent\makebox[0pt][r]{\smash{\raisebox{-15\baselineskip}{$(57)\ \left\{\rule{0pt}{30\baselineskip}\right.$}}\quad}a) Un groupe quotient $G_0$ du groupe cartographique triangulé $\unsure{\mathcal C_{\infty,3}|G_{\infty,3}|\mathfrak C_{\infty,3}}\simeq\{\tau_0,\tau_1,\tau_\infty\mid\tau_0^2=\tau_1^2=\tau_\infty^{2}=(\underbrace{\tau_1\tau_\infty}_{\rho_\infty})^3=(\underbrace{\tau_0\tau_\infty}_{\rho_1})^2=\}$

% L03–06
b) Un sous-groupe $\widetilde B''$ \deleted{\unsure{disons|du normalisateur}} \inserted{de $G_0$,} contenant le sous-groupe $\widetilde U$ engendré par $\tau_1,\tau_\infty$ (ou encore par $\rho_0=\tau_1\tau_\infty$, et $\tau_1$ ou $\tau_\infty$, $\widetilde U\simeq D_n$ \deleted{pour} où $n$ est l'ordre de $\rho_0$ dans $G_0$), normalisant $\widetilde U$

% L07–09
c) Un sous-groupe $B'$ de $\operatorname{Aut}_{\rm gr}(G_0\supset\widetilde B''\supset\widetilde U)$, \inserted{opérant trivialement sur $Z=\widetilde B''/\widetilde U$} contenant l'image de $\widetilde B''$ dans ce groupe par $g\mapsto\operatorname{int}_{G_0}(g)$, d'où un hom
\[
\psi:\widetilde B''\longrightarrow B'.
\]
% slanted
\margin{cela implique donc que $\widetilde B''/\widetilde U\simeq Z$ \unsure{commut.|commutatif}.}

% L10–16
[On suppose que $G_0$ opère fidèlement sur $G_0/\widetilde U$ (i.e l'intersection des conjugués de $\widetilde U$ dans $G_0$ est $\{1\}$) et que $B'$ opère fidèlement sur $G_0/\widetilde U$ par \unsure{transfert|transport} de structure -- la première condition impliquant déjà que $\psi$ induit un plongement
\[
\widetilde U\lhook\joinrel\xrightarrow{\psi\mid\widetilde U} B',
\]
on suppose $\widetilde U$ \uline{distingué} dans $B'$ \deleted{\unsure{d)}} d'où un hom déduit par passage au quotient
\[
Z\overset{\mathrm{def}}{=}\widetilde B''/\widetilde U\xrightarrow{\kappa}B'/\widetilde U\overset{\mathrm{déf}}{=}Z'.\quad\biggr]
\]

% L17–19
d) Un isomorphisme
\[
Z'\overset{\overline{\alpha}}{\underset{\vphantom{a}}{\overset{\sim}{\longleftrightarrow}}}Z
\]
rendant commutatif le diagramme
\diagram{\begin{tikzcd}[ampersand replacement=\&,column sep=large]
Z \arrow[r,"{\kappa}","{\text{cf c)}}"'] \arrow[dr,"{\lambda\mapsto\lambda^{2\deleted{\unsure{a}}}}"'] \& Z'\\
\& Z \arrow[u,"{\wr\ \alpha}"].
\end{tikzcd}}

% L20–26
Nous devons maintenant exprimer que \deleted{\unsure{des|dans}} la situation (26) correspondante à (57) \deleted{correspond} \inserted{provient} d'un graphe à homothéties, \deleted{\unsure{puis}} ayant \deleted{les arêtes \unsure{comme}} \inserted{$S$ comme} \deleted{$Z$} ens des sommets, ayant comme \deleted{\unsure{ens d'}} arêtes l'ens \unsure{arêtes|orienté} de la structure 2-polyédrale régulière correspondant à la donnée de $G_0$, et comme groupe d'homothéties $Z$; \deleted{Il \unsure{faut donc exprimer} axiomes \unsure{(1)}} \inserted{on suppose de plus, que $G$ soit le groupe des} autom. stricts du multigraphe à homothéties associé.

% archive sheet number: 175
